\documentclass[11pt, a4paper]{article}

\usepackage[margin=2.5cm]{geometry}
\usepackage[T1]{fontenc}
\usepackage[utf8]{inputenc}
% microtype disabled for compatibility
\usepackage{titlesec}
\usepackage{enumitem}
\usepackage{booktabs}
\usepackage{tabularx}
\usepackage{longtable}
\usepackage{array}
\usepackage{parskip}
\usepackage{fancyhdr}
\usepackage{xcolor}
\usepackage{mdframed}
\usepackage{listings}
\usepackage{hyperref}
\usepackage{amsmath}
\usepackage{tocloft}
\usepackage{multicol}

% Colors
\definecolor{rulecolor}{RGB}{40,40,40}
\definecolor{phasecolor}{RGB}{20,80,140}
\definecolor{warncolor}{RGB}{160,40,40}
\definecolor{codebackground}{RGB}{246,246,246}
\definecolor{codekeyword}{RGB}{0,80,160}
\definecolor{codestring}{RGB}{140,40,0}
\definecolor{codecomment}{RGB}{80,110,80}
\definecolor{tiercolor}{RGB}{60,120,60}

% Section formatting
\titleformat{\section}{\large\bfseries\color{phasecolor}}{}{0em}{}[\titlerule]
\titleformat{\subsection}{\normalsize\bfseries\color{phasecolor!70!black}}{}{0em}{}
\titleformat{\subsubsection}{\normalsize\itshape}{}{0em}{}

% Code listing style
\lstset{
  backgroundcolor=\color{codebackground},
  basicstyle=\ttfamily\footnotesize,
  keywordstyle=\color{codekeyword}\bfseries,
  stringstyle=\color{codestring},
  commentstyle=\color{codecomment}\itshape,
  breaklines=true,
  frame=single,
  framerule=0.4pt,
  rulecolor=\color{rulecolor!30},
  xleftmargin=8pt,
  xrightmargin=8pt,
  aboveskip=8pt,
  belowskip=8pt,
  showstringspaces=false,
  language=C++
}

% Header / footer
\pagestyle{fancy}
\fancyhf{}
\renewcommand{\headrulewidth}{0.4pt}
\fancyhead[L]{\small\texttt{UFX\_AUTO\_2} --- Full Design \& Work Order}
\fancyhead[R]{\small VSEPR-SIM / v5 beta8}
\fancyfoot[C]{\small\thepage}

% Rule box
\newmdenv[
  linecolor=warncolor,
  linewidth=1pt,
  leftmargin=0pt,
  rightmargin=0pt,
  innerleftmargin=10pt,
  innerrightmargin=10pt,
  innertopmargin=8pt,
  innerbottommargin=8pt,
  backgroundcolor=warncolor!4
]{rulebox}

% Phase / info box
\newmdenv[
  linecolor=phasecolor,
  linewidth=1pt,
  leftmargin=0pt,
  rightmargin=0pt,
  innerleftmargin=10pt,
  innerrightmargin=10pt,
  innertopmargin=8pt,
  innerbottommargin=8pt,
  backgroundcolor=phasecolor!4
]{phasebox}

% Tier box
\newmdenv[
  linecolor=tiercolor,
  linewidth=1pt,
  leftmargin=0pt,
  rightmargin=0pt,
  innerleftmargin=10pt,
  innerrightmargin=10pt,
  innertopmargin=8pt,
  innerbottommargin=8pt,
  backgroundcolor=tiercolor!4
]{tierbox}

% TOC tweaks
\renewcommand{\cftsecleader}{\cftdotfill{\cftdotsep}}
\setlength{\cftbeforesecskip}{2pt}

\begin{document}

% =========================================================
% TITLE
% =========================================================
\begin{center}
  {\LARGE\bfseries UFX\_AUTO\_2}\\[4pt]
  {\large\bfseries Full Design Document \& Work Order}\\[6pt]
  {\normalsize VSEPR-SIM $\cdot$ v5 beta8 $\rightarrow$ beta9 target}\\[3pt]
  {\small\texttt{Status: speculative architecture draft / not implementation gospel}}\\[3pt]
  {\small\textit{Because apparently we still believe in survival.}}
\end{center}

\vspace{4pt}
\hrule height 0.8pt
\vspace{10pt}

\tableofcontents
\newpage

% =========================================================
\section{Purpose}
% =========================================================

UFX\_AUTO\_2 expands the original UFX auto-creator from a UFF parameter completion tool into a broader materials-state database generator.

The first UFX auto system answered: \textit{Can we build, fill, validate, and log a UFF-style force-field table?}

UFX\_AUTO\_2 asks: \textit{Can we generate, validate, classify, and promote full material-state records across chemistry, thermodynamics, mechanics, transport, radiation, corrosion, and macro-scale engineering response?}

The goal is not merely to create more entries. The goal is a system that knows what kind of entry it is creating, what source class supports it, what uncertainty it carries, and whether it deserves to be promoted from generated to validated.

\begin{rulebox}
Tiny difference between ``database'' and ``organized hallucination farm.'' Best not to confuse them.
\end{rulebox}

\subsection{Core concept}

Every generated object is treated as a layered material record:

\begin{center}
\texttt{identity $\rightarrow$ atom typing $\rightarrow$ force field $\rightarrow$ molecular descriptors $\rightarrow$ thermodynamics $\rightarrow$ EOS/fluid $\rightarrow$ transport $\rightarrow$ mechanics $\rightarrow$ electronic $\rightarrow$ spectra $\rightarrow$ reaction/corrosion/radiation $\rightarrow$ macro observables $\rightarrow$ provenance $\rightarrow$ uncertainty $\rightarrow$ meta-scores}
\end{center}

The output is no longer only a \texttt{UFFEntry}. It becomes a \texttt{UFXMaterialRecord}.

\begin{lstlisting}
struct UFXMaterialRecord {
    IdentityBlock       identity;
    AtomTypeBlock       atom_types;
    ForceFieldBlock     force_field;
    MolecularDescriptorBlock molecular;
    CrystalBlock        crystal;
    ThermoBlock         thermo;
    EOSBlock            eos;
    TransportBlock      transport;
    MechanicalBlock     mechanical;
    ElectronicBlock     electronic;
    SpectralBlock       spectra;
    ReactionBlock       reactions;
    CorrosionBlock      corrosion;
    RadiationBlock      radiation;
    GranularBlock       granular;
    MacroObservableBlock macro;
    ProvenanceBlock     provenance;
    UncertaintyBlock    uncertainty;
    MetaScoreBlock      meta;
};
\end{lstlisting}

% =========================================================
\section{Core Identity Properties}
% =========================================================

These are the ``do not lose the atom'' fields. If identity is wrong, every downstream property is beautifully arranged garbage.

\subsection{Required fields}

\begin{multicols}{2}
\begin{itemize}[itemsep=1pt, leftmargin=*]
  \item element
  \item atomic number
  \item isotope
  \item mass
  \item formal charge
  \item oxidation state
  \item coordination number
  \item geometry tag
  \item hybridization / local environment
  \item atom type
  \item valence electron count
  \item electronegativity
  \item covalent radius
  \item ionic radius
  \item metallic radius
  \item van der Waals radius
  \item polarizability
\end{itemize}
\end{multicols}

\subsection{C++ struct}

\begin{lstlisting}
struct IdentityBlock {
    std::string material_key;
    std::string formula;
    std::string phase;

    std::vector<std::string> elements;
    std::vector<int> atomic_numbers;

    std::optional<int>    isotope_mass_number;
    std::optional<double> formal_charge;
    std::optional<double> oxidation_state;
    int coordination_number;
    std::string geometry_tag;
    std::string hybridization;

    std::string local_environment_hash;
};
\end{lstlisting}

\subsection{Uses}

Atom typing, fallback generation, bond prediction, charge balancing, coordination validation, periodic trends. PubChem PUG REST for identity cross-check. NIST WebBook for InChIKey and CAS.

% =========================================================
\section{UFF / Force-Field Properties}
% =========================================================

This is the traditional UFX core. Extended with audit and provenance fields so parameters stop wearing lab coats they did not earn.

\subsection{ForceFieldBlock fields}

\begin{multicols}{2}
\begin{itemize}[itemsep=1pt, leftmargin=*]
  \item $r_1$ --- bond radius
  \item $\theta_0$ --- preferred angle
  \item $x_1$ --- vdW distance
  \item $D_1$ --- vdW well depth
  \item $\zeta$ --- scaling / effective charge
  \item $Z_1$ --- effective charge parameter
  \item bond order correction
  \item coordination correction
  \item angle correction
  \item torsion barrier
  \item inversion barrier
  \item electrostatic charge model
  \item mixing rule metadata
\end{itemize}
\end{multicols}

\subsection{Extended audit fields}

\begin{itemize}[itemsep=1pt, leftmargin=*]
  \item \texttt{parameter\_source\_class}
  \item \texttt{parameter\_confidence}
  \item \texttt{parameter\_revision}
  \item \texttt{fallback\_rule}
  \item \texttt{local\_environment\_hash}
  \item \texttt{validated\_against}
\end{itemize}

\subsection{C++ struct}

\begin{lstlisting}
struct ForceFieldBlock {
    std::string atom_type;
    double r1, theta0, x1, D1, zeta, Z1;
    double bond_order_correction;
    double coordination_correction;
    double torsion_barrier;
    double inversion_barrier;
    std::string mixing_rule;
    std::string fallback_rule;
    std::string parameter_source_class;
    double parameter_confidence;
};
\end{lstlisting}

\subsection{Source classes}

\begin{itemize}[itemsep=1pt, leftmargin=*]
  \item \texttt{REFERENCE} --- Rapp\'{e} et al.\ 1992 UFF baseline
  \item \texttt{GENERATED} --- UFX\_AUTO\_2 fallback rules
  \item \texttt{VALIDATED} --- passed local, molecular, and web cross-checks
\end{itemize}

% =========================================================
\section{Molecular Descriptor Properties}
% =========================================================

Not deep physics. Excellent sanity filters. Store cheap descriptors on every molecule to catch the plausible catastrophes.

\subsection{Fields}

\begin{multicols}{2}
\begin{itemize}[itemsep=1pt, leftmargin=*]
  \item molecular formula
  \item molecular weight
  \item exact mass
  \item canonical SMILES
  \item InChI
  \item InChIKey
  \item formal charge
  \item heavy atom count
  \item rotatable bond count
  \item hydrogen bond donors
  \item hydrogen bond acceptors
  \item topological polar surface area
  \item XLogP
  \item complexity score
  \item fingerprint hash
\end{itemize}
\end{multicols}

\subsection{Failure modes these catch}

\begin{itemize}[itemsep=1pt, leftmargin=*]
  \item Generated \texttt{CH3HgI} entry claims wrong molecular weight
  \item Structure parser silently dropped iodine
  \item Charge balance is nonsense
\end{itemize}

Which, tragically, are all very plausible.

PubChem PUG REST exposes molecular properties and identifiers. Use it as an identity and descriptor cross-check source, not a force-field authority.

% =========================================================
\section{Thermochemical Properties}
% =========================================================

Required for gas, thermal, combustion, and phase behavior. State-dependent: do not store naked constants without attached state.

\subsection{Fields}

\begin{multicols}{2}
\begin{itemize}[itemsep=1pt, leftmargin=*]
  \item enthalpy of formation $\Delta H_f$
  \item Gibbs free energy $\Delta G_f$
  \item entropy $S$
  \item $C_p(T)$, $C_v(T)$
  \item enthalpy of combustion
  \item enthalpy of reaction
  \item equilibrium constants
  \item latent heat of fusion
  \item latent heat of vaporization
  \item boiling point
  \item melting point
  \item critical temperature $T_c$
  \item critical pressure $P_c$
  \item critical volume $V_c$
  \item acentric factor $\omega$
  \item vapor pressure curve
  \item Henry's law constant
\end{itemize}
\end{multicols}

\subsection{C++ struct}

\begin{lstlisting}
struct ThermoBlock {
    PropertyCurve cp_T;
    PropertyCurve h_T;
    PropertyCurve s_T;

    std::optional<double> melting_point_K;
    std::optional<double> boiling_point_K;
    std::optional<double> critical_temperature_K;
    std::optional<double> critical_pressure_Pa;
    std::optional<double> acentric_factor;
};
\end{lstlisting}

\subsection{Schema}

\begin{lstlisting}[language=SQL]
CREATE TABLE thermo_properties (
    id              INTEGER PRIMARY KEY AUTOINCREMENT,
    species_key     TEXT NOT NULL,
    property_name   TEXT NOT NULL,
    value           REAL,
    units           TEXT,
    temperature_K   REAL,
    pressure_Pa     REAL,
    source_id       TEXT,
    uncertainty     REAL,
    validation_status TEXT
);
\end{lstlisting}

NIST Chemistry WebBook is the primary reference for thermochemical data, phase transitions, vapor pressure, and reaction thermochemistry.

% =========================================================
\section{Equation-of-State and Fluid Properties}
% =========================================================

For gas/liquid/pipe/flow work. Store as functions, not single values. A single viscosity number with no temperature attached is a decorative lie.

\subsection{Fields}

\begin{multicols}{2}
\begin{itemize}[itemsep=1pt, leftmargin=*]
  \item density $\rho(T,P)$
  \item compressibility factor $Z(T,P)$
  \item specific volume
  \item dynamic viscosity $\mu(T,P)$
  \item kinematic viscosity $\nu(T,P)$
  \item thermal conductivity $k(T,P)$
  \item Prandtl number
  \item speed of sound
  \item bulk modulus
  \item surface tension
  \item diffusivity
  \item solubility
  \item fugacity coefficient
  \item activity coefficient
\end{itemize}
\end{multicols}

\subsection{Supported EOS models}

Ideal gas $\cdot$ van der Waals $\cdot$ Redlich-Kwong $\cdot$ Peng-Robinson $\cdot$ tabular steam/gas $\cdot$ mixed PvT tables

\subsection{C++ struct}

\begin{lstlisting}
struct EOSBlock {
    std::string eos_model; // ideal, VdW, RK, PR, tabular

    PropertySurface density_TP;
    PropertySurface Z_TP;
    PropertySurface viscosity_TP;
    PropertySurface thermal_conductivity_TP;
};
\end{lstlisting}

The property must be expressed as $f(T, P, \text{composition}, \text{phase})$.

% =========================================================
\section{Solid-State and Crystal Properties}
% =========================================================

Where UFX stops being molecule-only and starts touching materials science properly.

\subsection{Fields}

\begin{multicols}{2}
\begin{itemize}[itemsep=1pt, leftmargin=*]
  \item space group
  \item lattice parameters $a, b, c, \alpha, \beta, \gamma$
  \item unit cell volume
  \item density
  \item Wyckoff positions
  \item coordination polyhedra
  \item packing fraction
  \item nearest-neighbor distances
  \item formation energy
  \item energy above hull
  \item decomposition products
  \item phase stability
  \item magnetic ordering
  \item band gap
  \item density of states
  \item elastic tensor
  \item Debye temperature
  \item thermal expansion coefficient
\end{itemize}
\end{multicols}

\subsection{C++ struct}

\begin{lstlisting}
struct CrystalBlock {
    std::string space_group;
    Vec3 lattice_a_b_c;
    Vec3 alpha_beta_gamma_deg;

    double unit_cell_volume;
    double density;
    double formation_energy;
    double energy_above_hull;

    std::vector<Vec3> fractional_positions;
    std::string structure_source;
};
\end{lstlisting}

Materials Project and OQMD expose band gap, formation energy, magnetic ordering, elastic properties, and structural data for crystalline materials.

% =========================================================
\section{Mechanical Properties}
% =========================================================

For macro bridge and FEA coupling. The atomistic-to-engineering properties chain.

\subsection{Fields}

\begin{multicols}{2}
\begin{itemize}[itemsep=1pt, leftmargin=*]
  \item Young's modulus $E$
  \item shear modulus $G$
  \item bulk modulus $K$
  \item Poisson ratio $\nu$
  \item yield strength
  \item ultimate tensile strength
  \item fracture toughness $K_{IC}$
  \item hardness
  \item ductility
  \item strain hardening exponent
  \item fatigue strength
  \item creep constants
  \item stress relaxation constants
  \item plasticity model parameters
  \item damage evolution constants
  \item buckling load factor
\end{itemize}
\end{multicols}

\subsection{Scale chain}

\begin{center}
atomistic stiffness curvature $\rightarrow$ bead stiffness $\rightarrow$ field stress/strain $\rightarrow$ effective engineering constants
\end{center}

\subsection{C++ struct}

\begin{lstlisting}
struct MechanicalBlock {
    double E_eff;
    double G_eff;
    double K_eff;
    double poisson_ratio;

    double yield_strength;
    double failure_strength;
    double fracture_toughness;

    std::string test_method;
    double strain_rate;
    double temperature_K;
};
\end{lstlisting}

\subsection{Schema}

\begin{lstlisting}[language=SQL]
CREATE TABLE mechanical_properties (
    id           INTEGER PRIMARY KEY AUTOINCREMENT,
    material_key TEXT NOT NULL,
    test_type    TEXT NOT NULL,
    property_name TEXT NOT NULL,
    value        REAL,
    units        TEXT,
    strain_rate  REAL,
    temperature_K REAL,
    source_class TEXT,
    validation_score REAL
);
\end{lstlisting}

% =========================================================
\section{Transport Properties}
% =========================================================

The real bead-to-macro jackpot. Mandatory for pipeline erosion, reactive fluids, and radioactive materials.

\subsection{Fields}

\begin{multicols}{2}
\begin{itemize}[itemsep=1pt, leftmargin=*]
  \item thermal conductivity
  \item electrical conductivity
  \item ionic conductivity
  \item diffusion coefficient
  \item self-diffusion
  \item tracer diffusion
  \item species diffusivity
  \item permeability
  \item porosity
  \item tortuosity
  \item Darcy coefficient
  \item Forchheimer coefficient
  \item mass transfer coefficient
  \item heat transfer coefficient
  \item mobility
  \item Nernst-Einstein fields
\end{itemize}
\end{multicols}

\subsection{Macro connections}

\begin{itemize}[itemsep=1pt, leftmargin=*]
  \item bead trajectories $\rightarrow$ diffusivity
  \item particle-wall impacts $\rightarrow$ erosion flux
  \item thermal bead exchange $\rightarrow$ $k_\text{eff}$
  \item flow bead momentum exchange $\rightarrow$ pressure drop
\end{itemize}

\subsection{C++ struct}

\begin{lstlisting}
struct TransportBlock {
    double k_eff;
    double diffusivity;
    double permeability;
    double porosity;
    double tortuosity;
    double electrical_conductivity;
    double ionic_conductivity;
    double mass_transfer_coefficient;
};
\end{lstlisting}

% =========================================================
\section{Electronic and Quantum-Adjacent Properties}
% =========================================================

Not full quantum mechanics. A descriptor and constraint layer. Label the approximation before it humiliates you.

\subsection{Fields}

\begin{multicols}{2}
\begin{itemize}[itemsep=1pt, leftmargin=*]
  \item ionization energy
  \item electron affinity
  \item work function
  \item band gap
  \item HOMO-LUMO gap
  \item orbital occupations
  \item partial charges
  \item dipole moment
  \item quadrupole moment
  \item polarizability tensor
  \item dielectric constant
  \item effective mass
  \item density of states summary
  \item Fermi level
  \item redox potential
  \item electrochemical window
\end{itemize}
\end{multicols}

\subsection{UFX constraint examples}

\begin{itemize}[itemsep=1pt, leftmargin=*]
  \item High polarizability $\rightarrow$ stronger dispersion correction
  \item Low band gap $\rightarrow$ conductive/thermal coupling flag
  \item Large electron affinity $\rightarrow$ reactive capture / ionization risk flag
\end{itemize}

NIST WebBook includes ionization energies, electron affinities, proton affinities, and cluster ion binding energies. Materials Project includes electronic properties for crystalline materials.

% =========================================================
\section{Spectral and Fingerprint Properties}
% =========================================================

For identification and validation. Far better than ``the molecule didn't crash, so ship it.''

\subsection{Fields}

\begin{multicols}{2}
\begin{itemize}[itemsep=1pt, leftmargin=*]
  \item IR peaks
  \item Raman peaks
  \item UV/Vis absorption peaks
  \item mass spectrum peaks
  \item NMR shifts
  \item XRD peak list
  \item EXAFS / XANES signatures
  \item phonon spectrum
  \item vibrational frequencies
  \item normal modes
\end{itemize}
\end{multicols}

\subsection{UFX validation use}

\begin{center}
generated molecule $\rightarrow$ predicted vibrational modes $\rightarrow$ compare against IR/NIST peak ranges $\rightarrow$ score structural plausibility
\end{center}

NIST WebBook provides IR spectra, mass spectra, UV/Vis spectra, and vibrational constants for a large number of species.

% =========================================================
\section{Reaction and Kinetic Properties}
% =========================================================

How materials move from static properties to chemistry pathways. Where UFX becomes useful for pipelines, salts, reactors, and corrosive systems.

\subsection{Fields}

\begin{multicols}{2}
\begin{itemize}[itemsep=1pt, leftmargin=*]
  \item reaction enthalpy $\Delta H_\text{rxn}$
  \item reaction free energy $\Delta G_\text{rxn}$
  \item activation energy $E_a$
  \item transition state estimate
  \item Arrhenius $A$ factor
  \item rate constant $k(T)$
  \item reaction order
  \item equilibrium constant $K_\text{eq}$
  \item adsorption energy
  \item desorption energy
  \item surface binding site
  \item catalytic scaling descriptor
  \item reaction pathway ID
\end{itemize}
\end{multicols}

\subsection{C++ struct}

\begin{lstlisting}
struct ReactionBlock {
    std::string reaction_id;
    double delta_H;
    double delta_G;
    double activation_energy;
    double arrhenius_A;
    PropertyCurve rate_constant_T;
};
\end{lstlisting}

\subsection{System chain}

\begin{center}
solid + fluid + temperature + radiation $\rightarrow$ reactive surface $\rightarrow$ corrosion / adsorption / fouling / phase change
\end{center}

% =========================================================
\section{Radiation and Nuclear-Adjacent Properties}
% =========================================================

\begin{rulebox}
\textbf{Hard rule:} Do not store nuclear decay fields in the UFF parameter table. That table has suffered enough.
\end{rulebox}

\subsection{Fields}

\begin{multicols}{2}
\begin{itemize}[itemsep=1pt, leftmargin=*]
  \item isotope
  \item half-life
  \item decay mode
  \item decay energy
  \item alpha/beta/gamma probabilities
  \item neutron absorption cross section
  \item scattering cross section
  \item fission cross section
  \item activation product list
  \item radiolysis yield
  \item displacement threshold energy
  \item dpa rate proxy
  \item helium production rate
  \item gas swelling risk
  \item transmutation pathway
  \item decay heat
  \item shielding attenuation coefficient
\end{itemize}
\end{multicols}

\subsection{C++ struct}

\begin{lstlisting}
struct RadiationBlock {
    std::string isotope;
    double half_life_seconds;
    std::vector<std::string> decay_modes;
    double decay_heat_W_per_kg;
    double neutron_absorption_cross_section;
    double dpa_rate_proxy;
    double helium_production_rate;
    double swelling_risk;
};
\end{lstlisting}

\subsection{Scale chain}

\begin{center}
atom/bead identity $\rightarrow$ isotope-tagged material $\rightarrow$ decay event source term $\rightarrow$ heat/radiation damage field $\rightarrow$ macro swelling/failure/corrosion risk
\end{center}

Use separate tables: \texttt{nuclear\_properties}, \texttt{radiation\_response}, \texttt{decay\_channels}, \texttt{activation\_products}.

% =========================================================
\section{Corrosion and Environment Interaction Properties}
% =========================================================

For pipes, salts, molten systems, brines, acids, actinides.

\subsection{Fields}

\begin{multicols}{2}
\begin{itemize}[itemsep=1pt, leftmargin=*]
  \item oxidation tendency
  \item Pourbaix region
  \item corrosion potential
  \item pitting potential
  \item passivation tendency
  \item oxide layer growth rate
  \item solubility product
  \item chloride sensitivity
  \item fluoride sensitivity
  \item hydrogen embrittlement risk
  \item radiolysis coupling
  \item galvanic series relation
  \item salt compatibility
  \item redox buffer compatibility
\end{itemize}
\end{multicols}

\subsection{C++ struct}

\begin{lstlisting}
struct CorrosionBlock {
    std::string environment;
    double corrosion_potential;
    double pitting_risk;
    double passivation_score;
    double oxide_growth_rate;
    double salt_compatibility_score;
    double radiolysis_coupling_score;
    double hydrogen_embrittlement_risk;
};
\end{lstlisting}

\subsection{Macro connection}

\begin{center}
surface reaction $\rightarrow$ mass loss $\rightarrow$ wall thinning $\rightarrow$ stress concentration $\rightarrow$ failure risk
\end{center}

% =========================================================
\section{Powder, Granular, and Bead Properties}
% =========================================================

Directly supports DEM and pipeline erosion. This belongs in UFX macro expansion.

\subsection{Fields}

\begin{multicols}{2}
\begin{itemize}[itemsep=1pt, leftmargin=*]
  \item particle diameter
  \item shape factor
  \item sphericity
  \item aspect ratio
  \item roughness
  \item coefficient of restitution
  \item static friction
  \item rolling friction
  \item cohesion
  \item packing fraction
  \item coordination number
  \item angle of repose
  \item percolation threshold
  \item slugging risk
  \item segregation tendency
  \item abrasion coefficient
  \item erosion coefficient
\end{itemize}
\end{multicols}

\subsection{Pipeline erosion fields}

\begin{itemize}[itemsep=1pt, leftmargin=*]
  \item impact angle, impact velocity
  \item particle hardness, wall hardness
  \item restitution coefficient
  \item erosion yield, localized wear rate
\end{itemize}

\subsection{C++ struct}

\begin{lstlisting}
struct GranularBlock {
    double particle_diameter;
    double sphericity;
    double roughness;
    double restitution;
    double friction_static;
    double friction_rolling;
    double packing_fraction;
    double erosion_coefficient;
    double abrasion_coefficient;
};
\end{lstlisting}

% =========================================================
\section{Uncertainty and Provenance Properties}
% =========================================================

Not optional. This is what separates a database from a hallucination storage facility.

\subsection{Required fields on every property}

\begin{itemize}[itemsep=1pt, leftmargin=*]
  \item \texttt{source\_id} --- identifier of source record or URL
  \item \texttt{source\_class} --- see Section~2 enum
  \item \texttt{source\_version}
  \item \texttt{retrieval\_date}
  \item \texttt{raw\_response\_hash}
  \item \texttt{method\_tag}
  \item \texttt{confidence} --- 0.0--1.0
  \item \texttt{uncertainty} --- physical units
  \item \texttt{unit\_conversion\_trace}
  \item \texttt{validation\_status}
  \item \texttt{validation\_method}
  \item \texttt{number\_of\_confirmations}
  \item \texttt{conflict\_count}
  \item \texttt{last\_checked\_date}
  \item \texttt{promotion\_status}
\end{itemize}

\subsection{Schema}

\begin{lstlisting}[language=SQL]
CREATE TABLE property_provenance (
    id             INTEGER PRIMARY KEY AUTOINCREMENT,
    property_table TEXT NOT NULL,
    property_id    INTEGER NOT NULL,
    source_id      TEXT NOT NULL,
    source_class   TEXT NOT NULL,
    method_tag     TEXT,
    raw_hash       TEXT,
    confidence     REAL,
    uncertainty    REAL,
    retrieved_at   TEXT,
    notes          TEXT
);
\end{lstlisting}

\subsection{Provenance classification}

\begin{center}
\begin{tabular}{ll}
\toprule
\textbf{Status} & \textbf{Meaning} \\
\midrule
\texttt{published}   & from peer-reviewed or authoritative source \\
\texttt{DFT}         & from first-principles calculation \\
\texttt{generated}   & created by UFX\_AUTO\_2 rules \\
\texttt{inferred}    & derived from periodic trends \\
\texttt{cursed}      & should not be trusted \\
\bottomrule
\end{tabular}
\end{center}

That last category is wildly important.

% =========================================================
\section{Emergent UFX Meta-Properties}
% =========================================================

Not physical constants. System intelligence fields. These guide what to generate next.

\subsection{Fields}

\begin{itemize}[itemsep=2pt, leftmargin=*]
  \item \textbf{weirdness\_score} --- how unusual is this entry vs.\ periodic/structural neighbors?
  \item \textbf{coverage\_score} --- how well-represented is this region of parameter space?
  \item \textbf{confidence\_score} --- aggregate validation confidence
  \item \textbf{novelty\_score} --- distance from existing validated entries
  \item \textbf{database\_density\_score} --- local record density
  \item \textbf{local\_consistency\_score} --- agreement with immediate neighbors
  \item \textbf{cross\_source\_conflict\_score} --- inter-source disagreement
  \item \textbf{macro\_usefulness\_score} --- does this help predict $E_\text{eff}$, $k_\text{eff}$, $\Delta P$, damage, or phase stability?
  \item \textbf{simulation\_stability\_score} --- does this entry destabilize VSEPR runs?
  \item \textbf{extrapolation\_risk\_score} --- is the model outside validated space?
\end{itemize}

\subsection{Generation steering example}

\begin{phasebox}
\textbf{high novelty + low coverage + low extrapolation risk} $\rightarrow$ prioritize generation

\textbf{high weirdness + high conflict} $\rightarrow$ queue for manual review before promotion
\end{phasebox}

\subsection{C++ struct}

\begin{lstlisting}
struct MetaScoreBlock {
    double weirdness;
    double coverage;
    double confidence;
    double novelty;
    double local_consistency;
    double source_conflict;
    double macro_usefulness;
    double simulation_stability;
    double extrapolation_risk;
};
\end{lstlisting}

% =========================================================
\section{Full Far-Beyond Material Record Stack}
% =========================================================

The complete UFX object. Not one table. A layered record.

\subsection{Database structure}

\begin{center}
\begin{tabular}{ll}
\toprule
\textbf{Table} & \textbf{Owns} \\
\midrule
\texttt{material\_records}  & identity, class, timestamps \\
\texttt{property\_values}   & all property data, generic \\
\texttt{validation\_records} & comparisons and scores \\
\texttt{generation\_axes}   & randomized creation parameters \\
\texttt{promotion\_history} & class transitions \\
\texttt{property\_provenance} & source trail per property \\
\bottomrule
\end{tabular}
\end{center}

\subsection{Reconstruction query}

\begin{lstlisting}[language=SQL]
SELECT *
FROM material_records mr
JOIN property_values pv ON pv.material_id = mr.id
WHERE mr.material_key = ?;
\end{lstlisting}

\subsection{SourceClass enum}

\begin{lstlisting}
enum class SourceClass {
    Reference,    // trusted pinned source
    Generated,    // UFX_AUTO_2 rules or randomized
    Validated,    // passed local + external checks
    Rejected,     // failed sanity or validation
    Derived,      // computed from other stored quantities
    Imported,     // loaded from external database
    Experimental, // direct measured property
    Simulated     // produced by VSEPR/UFX simulation
};
\end{lstlisting}

Not all numbers deserve the same respect. Most numbers, like most committees, deserve suspicion.

% =========================================================
\section{Priority Implementation Order}
% =========================================================

Do not add everything at once unless the goal is a museum of half-working schemas.

\begin{tierbox}
\begin{tabular}{lll}
\toprule
\textbf{Tier} & \textbf{Blocks} & \textbf{Reason} \\
\midrule
0 & Identity + Provenance     & Know what the object is \\
1 & UFF / Force-field         & Know how it moves \\
2 & Molecular descriptors     & Sanity filters \\
3 & Thermo + EOS              & Know how state changes \\
4 & Crystal / material        & Know solid structure \\
5 & Transport + Mechanics     & Know how it flows and fails \\
6 & Reaction / Corrosion      & Know how it mutates \\
7 & Radiation / Nuclear       & Know how it degrades \\
8 & Spectra / Fingerprints    & Structural validation \\
9 & Meta-scores + Active learning & Know what to generate next \\
\bottomrule
\end{tabular}
\end{tierbox}

\begin{rulebox}
Do not implement Tier 7 before Tier 1 works. That is how people end up with a ``radiation-aware'' system that cannot parse oxygen. Majestic failure, but still failure.
\end{rulebox}

% =========================================================
\section{Randomized Generation Axes}
% =========================================================

UFX\_AUTO\_2 random generation samples a high-dimensional sweep. These axes drive the continual generation loop.

\subsection{Axis list}

\begin{multicols}{2}
\begin{itemize}[itemsep=1pt, leftmargin=*]
  \item element family
  \item isotope class
  \item oxidation state
  \item coordination
  \item geometry
  \item bond order
  \item temperature
  \item pressure
  \item phase
  \item density
  \item porosity
  \item crystal symmetry
  \item defect density
  \item dopant fraction
  \item surface orientation
  \item fluid environment
  \item radiation field
  \item corrosion medium
  \item particle size
  \item strain rate
  \item thermal gradient
  \item flow velocity
  \item validation source
  \item uncertainty target
\end{itemize}
\end{multicols}

\subsection{Axis config example}

\begin{lstlisting}[language={},basicstyle=\ttfamily\footnotesize]
{
  "element_family":  ["alkali","transition_metal","lanthanide",
                      "actinide","noble_gas"],
  "phase":           ["gas","liquid","solid","powder","molten_salt"],
  "coordination":    [1,2,3,4,5,6,8,12],
  "geometry":        ["linear","trigonal","tetrahedral",
                      "square_planar","octahedral"],
  "temperature_K":   [50, 2500],
  "pressure_atm":    [0.001, 1000],
  "defect_density":  [0.0, 0.10],
  "dopant_fraction": [0.0, 0.30],
  "strain_rate":     [1e-6, 1e3],
  "flow_velocity":   [0.0, 100.0]
}
\end{lstlisting}

\subsection{Example generated target}

\begin{lstlisting}[language={},basicstyle=\ttfamily\footnotesize]
{
  "element": "Fe",
  "oxidation_state": 2,
  "coordination": 6,
  "geometry": "octahedral",
  "phase": "solid",
  "defect_density": 0.015,
  "dopant": "Cr",
  "dopant_fraction": 0.12,
  "fluid_environment": "molten_fluoride_salt",
  "temperature_K": 973,
  "radiation_field": "gamma",
  "strain_rate": 0.001,
  "target_property": "corrosion_risk"
}
\end{lstlisting}

% =========================================================
\section{System Architecture and Generation Cycle}
% =========================================================

\subsection{System modules}

\begin{multicols}{2}
\begin{itemize}[itemsep=1pt, leftmargin=*]
  \item Identity Generator
  \item Atom-Type Generator
  \item Force-Field Generator
  \item Descriptor Importer
  \item Thermo/EOS Importer
  \item Crystal/Material Importer
  \item Mechanical Estimator
  \item Transport Estimator
  \item Reaction/Corrosion Estimator
  \item Radiation/Nuclear Tagger
  \item Granular/Bead Generator
  \item Web/Reference Validator
  \item SQLite Store
  \item Provenance Logger
  \item Uncertainty Scorer
  \item Meta-Score Engine
  \item Promotion Engine
\end{itemize}
\end{multicols}

\subsection{Generation cycle}

\begin{enumerate}[leftmargin=*, itemsep=2pt]
  \item Load coverage map
  \item Pick weak axis region
  \item Generate candidate material/state
  \item Fill identity block
  \item Fill force-field block
  \item Fill available descriptor/thermo/transport/mechanics blocks
  \item Run local sanity checks
  \item Store as \texttt{GENERATED}
  \item Queue for validation
  \item Validate from local / web / simulation references
  \item Promote, warn, or reject
  \item Update meta-scores
  \item Repeat
\end{enumerate}

\subsection{Cycle pseudocode}

\begin{lstlisting}
int run_auto2_cycle(const Auto2Options& opt) {
    SQLiteDB db(opt.db_path);
    CoverageMap coverage = db.load_coverage_map();
    AxisRegion target = coverage.choose_weak_region();

    UFXMaterialRecord candidate =
        generator.generate_material_record(target, opt.seed);

    LocalCheckResult local = local_checker.check(candidate);
    if (!local.pass) {
        db.insert_record(candidate, SourceClass::Rejected);
        return 1;
    }

    db.insert_record(candidate, SourceClass::Generated);

    ValidationBatch validation = validator.validate(candidate);
    double score = scorer.score(candidate, validation);

    if (score >= opt.promote_threshold)
        db.promote(candidate, SourceClass::Validated, score);
    else if (score >= opt.warn_threshold)
        db.mark_warn(candidate, score);
    else
        db.reject(candidate, score);

    db.update_meta_scores(candidate);
    return 0;
}
\end{lstlisting}

\subsection{Scoring model}

\begin{equation*}
S = 0.25 S_\text{identity} + 0.20 S_\text{local} + 0.20 S_\text{reference} + 0.15 S_\text{web} + 0.10 S_\text{repeat} + 0.10 S_\text{macro}
\end{equation*}

\begin{center}
\begin{tabular}{ll}
\toprule
\textbf{Score} & \textbf{Promotion class} \\
\midrule
$S \geq 0.95$          & \texttt{VALIDATED\_HIGH} \\
$0.85 \leq S < 0.95$   & \texttt{VALIDATED\_WARN} \\
$0.70 \leq S < 0.85$   & \texttt{GENERATED\_LOW\_CONFIDENCE} \\
$S < 0.70$             & \texttt{REJECTED} \\
\bottomrule
\end{tabular}
\end{center}

Missing webdata $\rightarrow$ confidence cap 0.85, not rejection. Missing provenance $\rightarrow$ automatic reject.

\subsection{CLI commands}

\begin{lstlisting}[language={}]
.\build\vsepr.exe ufx auto2 init
.\build\vsepr.exe ufx auto2 randomfill --count 500 --db ufx_auto2.sqlite
.\build\vsepr.exe ufx auto2 validate   --batch 100 --db ufx_auto2.sqlite
.\build\vsepr.exe ufx auto2 promote    --min-score 0.92 --db ufx_auto2.sqlite
.\build\vsepr.exe ufx auto2 audit      --db ufx_auto2.sqlite
.\build\vsepr.exe ufx auto2 run        --cycles 100 --batch 500 --web-batch 25
\end{lstlisting}

\subsection{Continual run harness}

Reuses existing PowerShell pattern from \texttt{uff\_continual}.

\begin{itemize}[leftmargin=*, itemsep=2pt]
  \item Run folders: \texttt{runs/ufx\_auto2/run\_YYYYMMDD\_HHMMSS\_NNN/}
  \item \texttt{summary.csv} per existing schema
  \item SQLite backup snapshot per cycle
  \item Web cache: \texttt{runs/ufx\_auto2/web\_cache/}
  \item Rate limiting on webdata fetcher
  \item \texttt{StopOnFail} and \texttt{CleanEachRun} flags
  \item Failure patterns: \texttt{missing=}, \texttt{nan}, \texttt{inf}, \texttt{fatal}, \texttt{exception}, \texttt{assert}
\end{itemize}

\subsection{Audit targets}

\begin{itemize}[leftmargin=*, itemsep=2pt]
  \item Record counts by class (Reference, Generated, Validated, Rejected)
  \item Coverage by element family, phase, property group, $T$/$P$ range
  \item Weakest axis region
  \item Highest conflict source
  \item Missing provenance count / missing validation count
  \item Highest weirdness entry / highest macro usefulness entry
\end{itemize}

\subsection{Risk register}

\begin{center}
\begin{tabular}{lll}
\toprule
\textbf{Risk} & \textbf{Problem} & \textbf{Mitigation} \\
\midrule
Source mixing     & Generated = published        & \texttt{source\_class} required everywhere \\
Webdata false fail & Missing = model failure     & \texttt{missing} status; confidence cap \\
Random spam       & Useless variants fill DB     & Coverage-aware scheduler + meta-scores \\
Schema bloat      & Premature dedicated tables   & \texttt{property\_values} first \\
Fake validation   & Local pass, no physics       & Multi-layer: local + ref + web + sim \\
\bottomrule
\end{tabular}
\end{center}

% =========================================================
\section{Hard Rules}
% =========================================================

\begin{rulebox}
\begin{enumerate}[leftmargin=*, itemsep=6pt]
  \item \textbf{No provenance = no promotion.} Period. Without it, UFX\_AUTO\_2 is a number blender with delusions of grandeur.
  \item \textbf{Nuclear/radiation fields never go in the UFF parameter table.} That table has suffered enough.
  \item \textbf{Missing webdata is not a failure.} It is a confidence cap of 0.85.
  \item \textbf{Generic \texttt{property\_values} table first.} Dedicated tables earn their place after usage stabilizes. Not before.
  \item \textbf{Do not implement Tier 7 (radiation) before Tier 1 (identity) is stable.} This has been said once. It will not be said again.
  \item \textbf{State is required.} A property value without attached $T$, $P$, and phase is a decorative lie.
  \item \textbf{Label the approximation.} Electronic and quantum-adjacent properties are constraint fields, not full QM. Say so.
  \item \textbf{SQLite is truth. JSONL is history. CSV is convenience. Binary is disposable.} Treat them accordingly.
\end{enumerate}
\end{rulebox}

\vspace{10pt}
\hrule
\vspace{8pt}

\begin{center}
\textit{The database should know what it knows, what it guessed, and what it made up.\\
All three categories exist. Only the last one is embarrassing.}
\end{center}

\end{document}
